\section{人工智能与机器学习}

\subsection{制造业里的 AI 与 ML}

理解机器学习之前，有必要先区分两个常被混用的概念。强人工智能主张 AI 系统可以具备和人类相同、甚至超越人类的智力能力；弱人工智能则聚焦于用数学和信息学方法解决具体的应用问题，系统具备一定的自我优化能力，过程中会模拟人类智能的某些方面，或者构建用于模拟、辅助人类思考的系统。目前工业界能落地的绝大多数应用，属于弱人工智能的范畴。

而机器学习的核心机制是：机器从已有的样本数据中识别模式，据此构建模型，再把学到的知识应用到此前未见过的新情况上。数据量越大、信息越丰富，模型学到的效果通常越好。

生产领域被认为是机器学习应用潜力最大的场景之一，服务、客户服务和物流等领域同样存在丰富的应用空间。机器学习在系统自我优化与自我适应方面展现出很强的潜力。

几个因素正在加速机器学习技术在工业场景中的落地：开发环境的显著进步降低了使用门槛；特征工程这一原本繁琐的环节被大大简化；性能优异的预训练网络结构（尤其在图像分析与识别领域）可以直接复用；专用且成本更低的硬件也加快了模型训练的速度。这几个使能因素共同解释了为什么近年来机器学习方法能够更快地走出实验室、进入实际产线。

\subsection{图像特征}

机器视觉要解决的根本问题可以归结为两个问题：人类究竟依据哪些信息（图像特征）在视觉上识别出物体？这些信息又该如何表示，才能让计算机也具备识别物体的能力？

从算法角度看，这个问题可以进一步拆成两个子任务。分类任务是把一次观测分配给某个类别；而机器学习的任务则是学习出一组类别边界，使这种分配尽可能不出错。

对图像特征本身，有三点基本要求：可靠性（同一物体在不同条件下能被稳定识别）、可计算性（提取特征所需的计算开销要可控）、可区分性（不同类别的特征要能被明确区分开）。图像特征大致可以分成两类：一类是特定的强度变化（亮度变化），可以通过梯度大小、梯度方向、卷积操作（滤波器）等方式量化计算；另一类是图像局部区域内强度变化在特定方向上的分布，常用方向直方图来描述。

在卷积神经网络普及之前，图像特征提取主要依赖两条技术路线：基于梯度的方法，以及基于统计分布的方法。

\paragraph{基于梯度的方法}

Moravec 角点检测器是较早的思路：取图像中一个小窗口，把这个窗口分别向不同方向做微小平移，观察平均亮度的变化情况——几乎没有变化说明处在平坦区域；只在特定方向平移时亮度发生明显变化说明处在边缘上；无论向哪个方向平移亮度都明显变化则说明处在角点上。这种方法的缺点是滤波响应具有各向异性，而且二值化的矩形窗口容易受噪声干扰。

Harris 角点检测器针对这个问题做了改进：通过每个像素点处灰度函数的偏导数构建一个结构矩阵，矩阵的特征值反映边缘强度，特征向量反映边缘方向，特征值可以通过矩阵的迹和行列式求出。这样得到的滤波响应具有各向同性，对平移和旋转具有不变性，对亮度变化也有一定的鲁棒性，但缺点是不具备尺度不变性。

Haar 特征最初是为人脸检测设计的，思路来自 Haar 小波（与傅里叶分析类似）：对图像区域内的强度做乘法和求和运算，强度相近时得到较低的数值，强度差异大时得到较高的正负数值；借助积分图技术可以高效计算像素求和，大幅提升计算速度。

\paragraph{基于统计分布的方法}

方向梯度直方图（Histogram of Oriented Gradients，HOG）先做预处理，再对每个像素通过卷积运算计算梯度的大小和方向；接着把图像窗口划分成若干局部单元格，每个单元格内的梯度汇总成一个方向直方图；最后对重叠的区块做归一化处理。HOG 的优点是得到的特征向量可以直接作为分类器的输入，且计算速度快、特征表达能力强；缺点是需要预先设定固定的窗口尺寸，对旋转、尺度和平移变化的鲁棒性有限。

尺度不变特征变换（Scale Invariant Feature Transform，SIFT）则是为了解决尺度和旋转不变性的问题：它能找到图像中的关键点（Keypoints）并用特征向量描述，这种描述对尺度和旋转具有不变性，同时对仿射畸变、视角变化、噪声和光照变化都有较好的鲁棒性。SIFT 算法大致分四步：先在尺度空间中确定候选关键点，尺度空间的构建依赖对同一图像做不同程度的高斯模糊，再把相邻模糊程度的图像相减（高斯差分），这是一种能快速突出边缘、生成梯度图像的方法；接着筛选和定位关键点；再为每个关键点分配一个主方向，把邻域按主方向旋转以获得旋转不变性；最后用邻域内的方向直方图描述每个关键点，得到一个 128 维的特征向量。

\subsection{卷积神经网络}

经典方法有一个共同的痛点：特征提取环节高度依赖人工设计和专家经验，需要反复尝试和评估不同的特征提取与分类方案，费时费力。卷积神经网络（CNN）最大的价值就在于取消了这道特征工程的门槛。特征提取被整合进分类模型内部，训练过程中自动完成，这种从原始图像直接到分类结果的方式被称为端到端学习（End-to-End Learning）。

如图 \ref{fig:cnn} 所示：CNN 的核心组件是卷积层：神经元按宽度、高度、深度三个维度排列，专门为图像处理设计。卷积层的参数由一组可学习的感受野滤波器（权重矩阵）组成，这些滤波器也被称为特征检测器，其宽高通常远小于输入图像，滤波器的数量对应输入图像的深度。在每个位置上，滤波器的各个通道与图像对应区域的通道逐元素相乘再求和，这个操作就是卷积；一个滤波器在整张图像上卷积得到的结果构成一张特征图（Feature Map），特征图中被突出显示的区域，就是与该滤波器最相似的区域。

\begin{figure}[htbp]
  \centering
  \includegraphics[width=0.8\textwidth]{cnn.jpg}
  \caption{卷积神经网络}
  \label{fig:cnn}
\end{figure}

池化层通常夹在两个卷积层之间，作用是从前一卷积层得到的特征图中剔除对模式识别而言冗余的参数，从而降低计算量。通过多层卷积和池化的堆叠，网络可以构建出复杂的非线性映射关系：层与层之间形成一种层级结构，逐步实现自动特征提取，特征也随层数加深逐级浓缩。每个节点基于前一层的输出学习特定的特征——最初的几层通常只识别角点、边缘等简单模式，越往后的层则负责聚合和组合这些简单模式，从而识别出更复杂的形状和物体。

\subsection{图像内容分类}

经典分类方法中，贝叶斯分类器基于贝叶斯定理
\begin{equation}
    P(c|x) = \frac{P(x|c)\cdot P(c)}{P(x)}
\end{equation}
对给定的观测 \(x\)，选择使后验概率 \(P(c|x)\) 最大的类别作为判断结果。它的优点是模型简单、易于实现，且所需的训练数据量较少。

k 近邻（k-Nearest-Neighbor）走的是一条更直观的路线，如图 \ref{fig:knn} 所示：判断一个未知物体属于哪一类，依据是它和已知物体之间的相似度，相似度通过距离度量转换成具体数值来表达。优点是分类器直观、实现简单，同时又能够拟合比较复杂的类别边界。

\begin{figure}[htbp]
  \centering
  \includegraphics[width=0.4\textwidth]{knn.png}
  \caption{k 近邻分类器}
  \label{fig:knn}
\end{figure}

如图 \ref{fig:svm} 所示，支持向量机（SVM）的核心思路是用一个超平面把线性可分的数据分开，并让训练样本（支持向量）到分界面的正交距离尽可能大；对于线性不可分的数据，可以通过映射到更高维的特征空间来处理，这就是常说的核技巧（Kernel-Trick）。

\begin{figure}[htbp]
  \centering
  \includegraphics[width=0.4\textwidth]{svm.png}
  \caption{支持向量机}
  \label{fig:svm}
\end{figure}

神经网络的基本结构是：输入值通常写成向量 \(x=(x_1, x_2, \ldots, x_m)\)，每个输入对应一个网络权重 \(w_i\)，网络活性 net 由输入与权重的加权和计算得出，再经过一个激活函数 \(f(x)\) 得到输出。可以调节的超参数包括网络架构或拓扑结构（神经元数量、层数、连接方式）、激活函数的类型（阈值函数、Sigmoid 函数等）、训练过程中权重的调整方式（反向传播算法），以及阈值的调整方式。

用CNN做分类时，网络末端通常接若干全连接层（Fully Connected Layers，FC）。倒数第二个全连接层输出一个长度为 \(k\) 的向量，\(k\) 等于类别数量；对这一层应用 Softmax 激活函数，计算出每个类别对应的概率，所有类别的概率加起来正好等于 1。